\documentclass[11pt,a4paper]{article}

\usepackage[utf8]{inputenc}
\usepackage[T1]{fontenc}
\usepackage[french]{babel}
\usepackage[margin=2.5cm]{geometry}
\usepackage{booktabs}
\usepackage{amsmath}
\usepackage{verbatim}
\usepackage[hidelinks]{hyperref}

\setlength{\parindent}{0pt}
\setlength{\parskip}{6pt}

\title{TP : Paralléliser un calcul de moyenne avec Numba}
\author{Ibrahima Gabar Diop \\ \small M1 MBDA, Université Cheikh Hamidou Kane}
\date{Juin 2026}

\begin{document}
\maketitle

\section{Objectif}

Calculer la moyenne d'une promo de $N$ étudiants, d'abord en séquentiel puis en
parallèle avec Numba, mesurer le gain de temps (speedup) et estimer la part
parallélisable du calcul avec la loi d'Amdahl.

\section{Le calcul}

Chaque étudiant a trois notes pondérées par des coefficients.

\begin{center}
\begin{tabular}{lc}
\toprule
Matière & Coefficient \\
\midrule
Maths (M) & 5 \\
Physique (P) & 4 \\
Anglais (A) & 2 \\
\bottomrule
\end{tabular}
\end{center}

La moyenne d'un étudiant vaut $(M \times 5 + P \times 4 + A \times 2) / 11$. On la
calcule pour chacun, puis on en fait la moyenne sur toute la promo. Les étudiants
sont indépendants : aucune note ne dépend d'une autre, donc la boucle peut être
répartie sur plusieurs threads sans fausser le résultat.

\section{Méthode}

Les données sont générées dans un fichier CSV (notes tirées entre 0 et 20). Le
calcul est écrit en deux versions, toutes deux compilées par Numba.

La version séquentielle utilise le décorateur \texttt{@njit} et s'exécute sur un
seul thread : elle sert de référence. La version parallèle utilise
\texttt{@njit(parallel=True)} avec \texttt{prange} : la boucle est répartie sur
les threads et la somme finale est gérée comme une réduction, sans condition de
course.

Le coeur de la version parallèle :

\begin{verbatim}
@njit(parallel=True, cache=True)
def moyenne_promo_par(M, P, A):
    total = 0.0
    n = M.shape[0]
    for i in prange(n):
        total += (M[i] * 5 + P[i] * 4 + A[i] * 2) / 11
    return total / n
\end{verbatim}

\section{Résultats}

Mesure sur 10 000 000 d'étudiants, machine à 4 coeurs physiques (8 threads
logiques). Le temps retenu est le meilleur sur 100 exécutions, compilation JIT
exclue. Les deux versions donnent exactement le même résultat (moyenne de la
promo : $10{,}0004$).

\begin{center}
\begin{tabular}{ccc}
\toprule
Threads & Temps (ms) & Speedup \\
\midrule
Séquentiel & 18,28 & 1,00 \\
1 & 17,20 & 1,06 \\
2 & 9,35 & 1,96 \\
4 & 8,66 & 2,11 \\
8 & 9,88 & 1,85 \\
\bottomrule
\end{tabular}
\end{center}

Avec un seul thread, la version parallèle rejoint la séquentielle : le surcoût de
\texttt{prange} est négligeable. Le speedup est presque idéal à 2 threads (1,96)
et atteint son maximum à 4 threads (2,11), soit le nombre de coeurs physiques. À
8 threads il redescend, car les threads supplémentaires sont des hyperthreads qui
partagent les mêmes coeurs.

\section{Speedup et loi d'Amdahl}

Le meilleur speedup mesuré est $S = 2{,}11$ pour $n = 4$ threads. La loi d'Amdahl
relie le speedup à la proportion parallélisable $p$ :

\[
S = \frac{1}{(1 - p) + \dfrac{p}{n}}
\qquad\Longrightarrow\qquad
p = \frac{1 - \dfrac{1}{S}}{1 - \dfrac{1}{n}}
\]

En remplaçant :

\[
p = \frac{1 - \dfrac{1}{2{,}11}}{1 - \dfrac{1}{4}}
  = \frac{0{,}526}{0{,}75}
  \approx 0{,}70
\]

La proportion parallélisable du calcul est donc estimée à environ 70 \%.

\section{Discussion}

Le calcul fait peu d'opérations par étudiant mais lit beaucoup de données : il
est en partie limité par la bande passante mémoire. C'est pourquoi le speedup
plafonne autour de 2 au lieu d'atteindre le nombre de coeurs, et pourquoi
l'hyperthreading n'apporte rien. La tâche est parallèle à 100 \% sur le papier,
mais limitée en pratique par l'accès mémoire, ce que la loi d'Amdahl traduit par
un $p$ inférieur à 1. Les valeurs exactes varient un peu selon la charge de la
machine ; l'ordre de grandeur (speedup proche de 2, $p$ entre 60 et 70 \%) reste
stable.

\section{Reproduire}

\begin{verbatim}
pip install -r requirements.txt
python generate_data.py --n 10000000 --out etudiants.csv
python benchmark.py --csv etudiants.csv --repeat 100
\end{verbatim}

Code complet : \url{https://github.com/Gblack98/tp-parallelisation-numba}

\end{document}
